\begin{abstract}
Estimating the evolution of the battery's State of Charge (SoC) in response to its usage is critical for implementing effective power management policies and for ultimately improving the system's lifetime.
Most existing estimation methods are either physics-based digital twins of the battery or data-driven models such as Neural Networks (NNs).
In this work, we propose two new contributions in this domain. First, we introduce a novel NN architecture formed by two cascaded branches: one to predict the \textit{current} SoC based on sensor readings, and one to estimate the SoC at a \textit{future} time as a function of the load behavior. Second, we integrate battery dynamics equations into the training of our NN, merging the physics-based and data-driven approaches, to improve the models' generalization over variable prediction horizons. 
We validate our approach on two publicly accessible datasets, showing that our Physics-Informed Neural Networks (PINNs) outperform purely data-driven ones while also obtaining superior prediction accuracy with a smaller architecture with respect to the state-of-the-art.  

% \rev{
% \st- \st{cambiamento Tab 1 per aggiungere mem e ops nostri e di SoA;}

% \st{- chiarire metodi che ground truth è usata solo per training;}

% - \st{aggiungere espeirmenti Sandia con tutte temperautre / cicli di scarica;}

% - \st{CONFRONTARSI CON "UNICO?" PAPER CHE FA PREDICTION; FARE RIASSUNTO E DISCUTERNE; CONTROLLARE ANCHE ALTRI PAPERS}

% \st{- sottolineare che non consideriamo SOH e il nostro algoritmo è ortogonale: puoi trainare SOH model e un SoC prediction a ogni diverso livello di SOH}
% }
\end{abstract}
